\documentclass{beamer}

% === AUTOR === (((
\author{\textit{Por Erick I. Rodríguez Juárez.}}
% )))

% === PAQUETES === (((
% \usepackage{makeidx}
% \usepackage{xltxtra}
\usepackage{amsfonts}
\usepackage{amsmath}
\usepackage{amssymb}
% \usepackage{fullpage}
\usepackage{tikz}
\usetikzlibrary{arrows.meta}
\usepackage{graphicx}
% )))

% === TIPOGRAFÍA === (((
% \setmainfont[
  % BoldFont       = bodonibi,
	% ItalicFont     = Century modern italic2.ttf,
	% BoldItalicFont = bodonibi,
	% SmallCapsFont  = lmromancaps10-regular.otf
% ]{Century_modern.ttf}
% )))

% === COMANDOS === (((
% \newcommand{\dis}{\displaystyle}
% \newcommand{\qed}{\hspace{0.5cm}\rule{0.16cm}{0.4cm}}
% \newcommand{\operator}[1]{\mathop{\vphantom{\sum}\mathchoice
% {\vcenter{\hbox{\huge $#1$}}}
% {\vcenter{\hbox{\Large $#1$}}}{#1}{#1}}\displaylimits}
% \newcommand{\suma}{\operator{\includegraphics[scale=0.09]{FOTOS/Sigma.png}}}
% \setlength{\parindent}{0mm}
% )))

% === ITALICA EN ENTORNO MATEMÁTICO === (((
% \DeclareSymbolFont{italics}{\encodingdefault}{\rmdefault}{m}{it}
% \DeclareSymbolFontAlphabet{\mathit}{italics}
% \ExplSyntaxOn
% \int_step_inline:nnnn { `A } { 1 } { `Z }
 % {  \exp_args:Nf \DeclareMathSymbol{\char_generate:nn{#1}{11}}{\mathalpha}{italics}{#1} }
% \int_step_inline:nnnn { `a } { 1 } { `z } {  \exp_args:Nf \DeclareMathSymbol{\char_generate:nn{#1}{11}}{\mathalpha}{italics}{#1}}
% \ExplSyntaxOff
% )))

\begin{document}

\frame{\titlepage}

\begin{frame}[t]
	\frametitle{Teoría Básica para Ecuaciones Diferenciales Lineales Homogéneas.}
\begin{definition}
	Una E.D. Lineal de Orden \(n\) de la forma
	\[
		a_n(x) y^{(n)} +a_{n-1} (x) y^{(n-1)} + \;\cdots\; + a_1(x) y' +a_0(x) y = g(x).
	\]
	se conoce como \textbf{homogénea} si \(g(x) \equiv 0\), mientras que si \(g(x)\) no es idénticamente cero, se conoce como \textbf{no homogénea}.
\end{definition}
	\begin{example}
		\[
			\begin{array}{rcl}
				3x^2y''' + \dfrac{1}{x} y - \sin x & = & 0 \\[2mm]
				3x^2y''' + \dfrac{1}{x} y & = & \underbrace{\sin x} _{g(x) \ne 0} .
			\end{array}
		\]
		Por tanto, es una E.D. no homogénea.
	\end{example}
\end{frame}

\begin{frame}[t]
	\begin{example}
		\[
			y^{(iv)} +y'' = \underbrace{0} _{g(x) =0}.
		\]
		Por tanto, es una E.D. homogénea.
	\end{example} 
	\begin{block}{Principio de Superposición.}
		Sean \(y_1,y_2, \;\cdots\; ,y_k\) soluciones de la E.D. lineal homogénea de orden \(n\) en el intervalo \(I\). Entonces, la combinación lineal
		\[
			y(x) = c_1y_1(x) + c_2y_2(x) + \;\cdots\; c_ky_k(x).
		\]
		Es solución de la E.D. homogénea, donde \(c_1, \;\ldots,\; c_k\) son constantes arbitrarias.
	\end{block}
\end{frame}

\begin{frame}[t]
	\begin{corollary}
		\begin{enumerate}
			\item Un múltiplo constante de \(y=cy_1(x)\) de una solución \(y_1(x)\) de una E.D. lineal homogénea también es solución.
			\item Una E.D. lineal homogénea siempre tiene la solución trivial \(y \equiv 0\).
		\end{enumerate}
	\end{corollary} 
	\begin{example}
		Las funciones \(y_1(x) = x^2\), \(y_2(x) =x^2 \ln x\) son soluciones de la ecuación diferencial lineal homogénea \(x^3y''' -2xy' +4y=0\), en \((0, \infty)\). Determina un número infinito de soluciones usando el principio de superposición.
	\end{example}
	\begin{exampleblock}{}
		\textbf{Solución.} Usando el principio de superposición \(y(x) = c_1x^2+c_2x^2 \ln x\), \(c_1,c_2\) constantes arbitrarias. Con ésta expresión se tiene un número infinito de soluciones de la ecuación diferencial. \\[2mm]
	\end{exampleblock}
\end{frame}

\begin{frame}[t]
	\begin{example}
		Si \(y(x) = e^{7x}\) es solución de la ecuación diferencial \(y'' -9y' +14y=0\), encuentra 3 soluciones de la ecuación diferencial.
	\end{example}
	\begin{exampleblock}{}
		\textbf{Solución.} 
		\begin{enumerate}
			\item \(y_1(x) = 3e^{7x}\).
			\item \(y_2(x) = 0\).
			\item \(y_3(x) = - \sqrt{2} e^{7x}\).
		\end{enumerate}
	\end{exampleblock}
\end{frame}

\begin{frame}[t]
	\begin{block}{Independencia Lineal en términos del Wronskiano.}
		Se dice que un conjunto de funciones \(f_1(x) , \;\cdots\; f_n(x)\) es \textbf{linealmente dependiente} (LD) en un intervalo \(I\) si existen constantes \(c_1, \;\ldots,\; c_n\) no todos ceros tales que
		\[
			c_1f_1(x) + c_2f_2(x) + \;\cdots\; c_nf_n(x) =0, \hspace{5mm} \mbox{para todo } x \in I.
		\]
		Si el conjunto no es LD en el intervalo \(I\), se considera que es \textbf{linealmente independiente} (LI), (es decir, cuando las únicas constantes para que se cumpla que la combinación lineal sea cero, son \(c_1=c_2= \;\cdots\; c_n=0\)).
	\end{block} 
	En particular, nos interesa determinar cuando son LI y eso puede determinarse mecánicamente recurriendo a un determinante llamado \textbf{Wronskaino}.
\end{frame}

\begin{frame}[t]
	\begin{block}{Definición. (Wronskiano)}
		Suponga que cada una de las funciones \(f_1(x) , \;\cdots\; f_n(x)\) posee al menos \(n-1\) derivadas. Luego a
		\[
			W(f_1, \;\cdots\; f_n) = \left| \begin{array}{llll}
				f_1 				& f_2 				& \;\cdots\;  				& f_n\\[2mm]
				f_1' 				& f_2' 				& \;\cdots\;  				& f_n'\\[2mm]
				\vdots & \vdots & \ddots & \vdots \\[2mm]
				f^{(n-1)}_1 				& f^{(n-1)}_2 				& \;\cdots\;  				& f^{(n-1)}_n
			\end{array}\right| 
		\]
		se le conoce como \textbf{Wronskiano} de las \(n-\)funciones.
	\end{block} 
	\begin{block}{Criterio para Determinar Soluciones LI.}
		Sean \(y_1, \;\ldots,\; y_n\), \(n\) soluciones de la ecuación diferencial lineal homogénea de orden \(n\) en un intervalo \(I\). Entonces el conjunto de soluciones es LI, en \(I\), si y sólo si \(W(y_1, \;\ldots,\; y_n) \ne 0\) para todo \(x \in I\).
	\end{block}
\end{frame}

\begin{frame}[t]
	\begin{block}{Conjunto Fundamental de Soluciones.}
		\textbf{Definición.} Es un conjunto de \(n\) soluciones LI de la ecuación diferencial homogénea de orden \(n\) en un intervalo \(I\).
	\end{block}
	\begin{block}{Solución General de una E.D. Homogénea.}
		\textbf{Definición.} Sean \(y_1, \;\ldots,\; y_n\) un conjunto de fundamental de soluciones de la E.D. homogénea de orden \(n\), en un intervalo \(I\). Entonces la solución general de la ecuación en el intervalo es:
		\[
			y(x) = c_1y_i(x) + \;\cdots\; +c_ny_n(x).
		\]
		donde \(c_1, \;\ldots,\; c_n\) son constantes positivas.
	\end{block}
\end{frame}

\begin{frame}[t]
	\begin{block}{Procedimiento para Resolver E.D. Lineales Homogéneas.}
		Para determinar todas las soluciones dela E.D.
		\[
			a_n(x) y^{(n)} + \;\cdots\; + a_1(x) y' +a_0(x) y =0,
		\]
		\vspace{-6mm}
		\begin{enumerate}
			\item Determina \(n\) soluciones \(y_1, \;\ldots,\; y_n\) que constituyan un conjunto fundamenta de soluciones.
			\item Dar la solución general de la ecuación diferencial como:
			\[
				y(x) = c_1y_1(x) + \;\cdots\; +c_ny_n(x).
			\]
		\end{enumerate}
	\end{block}
	\begin{example}
		Sean \(y_1(x) = e^{x/2}\), \(y_2(x) = xe^{x/2}\) soluciones de la E.D.L.H. \(y'' -y' +1/4y=0\),  en el intervalo de \((- \infty , \infty)\). ¿Es posible dar la solución general de la E.D.? Si es así, escríbala explícitamente.
	\end{example}
\end{frame}

\begin{frame}[t]
	\begin{exampleblock}{}
		\textbf{Solución.} Veamos si \(y_1\) y \(y_2\) son L.I. usando el Wronskiano.
		\[
			\begin{array}{rl}
				W[y_1,y_2] (x) & = \begin{vmatrix}
					e^{x/2} & xe^{x/2} \\
					1/2e^{x/2} & 1/2xe^{x/2} +e^{x/2}
				\end{vmatrix}\\[6mm]
				& = e^{x/2} (1/2xe^{x/2} +e^{x/2}) -1/2e^{x/2} (xe^{x/2}) \\[2mm]
				& = 1/2xe^x+e^x-1/2xe^x = e^x \ne 0, \forall x \in (- \infty , \infty) .
			\end{array}
		\]
		\(\therefore \hspace{5mm} y_1\) y \(y_2\) son linealmente independientes (L.I.) Luego \(\{y_1,y_2\}\) conforman un conjunto fundamental de soluciones y por consiguiente la solución general de la E.D. está dada por
		\[
			y(x) = c_1e^{x/2} + c_2xe^{x/2} , \hspace{5mm} c_1,c_2 \mbox{ constantes arbitrarias.}
		\]
	\end{exampleblock}
\end{frame}

\begin{frame}[t]
	\begin{alertblock}{Ejercicio.}
		Sea \(y_2(x) = e^x \cos 2x\) , \(y_2(x) = e^x \sin 2x\) soluciones de la E.D.L.H. \(y'' +2y' +5y=0\) en el intervalo \((- \infty , \infty)\).
		Compruebe que el conjunto que contiene a \(y_1,y_2\) es un conjunto fundamental de soluciones y escriba explícitamente la solución general de la E.D.\\
		\textbf{Solución.} Veamos si \(y_1\) y \(y_2\) son L.I.  \vspace{-3mm}
		\[
			\begin{array}{rl}
				W[y_1,y_2] (x) = & \begin{vmatrix}
					e^x \cos 2x & e^x \sin 2x \\
					-2e^x \sin 2x+e^x \cos 2x & 2e^x \cos 2x+e^x \sin 2x
				\end{vmatrix}\\[2mm]
				 = & e^x \cos 2x(2e^x \cos 2x+e^x \sin 2x) - \\[2mm]
				& - (-2e^x \sin 2x+ e^x \cos 2x) e^x \sin 2x \\[2mm]
				 = & 2x^{2x} \cos ^22x+ \cancel{e^{2x} \cos 2x \sin 2x} +2e^{2x} \sin ^2 2x  \\[2mm]
				& - \cancel{e^{2x} \cos 2x \sin 2x} \\[2mm]
				 = & 2e^{2x} \cos ^22x +2e^{2x} \sin ^22x \\
				 = & 2e^{2x} \cancelto{1}{(\cos ^22x+ \sin ^22x)} = 2e^{2x} \ne 0.
			\end{array} \vspace{-4mm}
		\]
		\(e^{2x} \ne 0\), \(\forall x \in (- \infty , \infty)\).
	\end{alertblock}
\end{frame}

\begin{frame}[t]
	\begin{alertblock}{}
		\begin{center}
			\(\therefore \hspace{5mm} y_1 \;,\; y_2\) son L.I.
		\end{center}
		Dado que la E.D. es de \(2^{do}\) orden, entonces \(\{y_1,y_2\}\) conforman una C.F.S y por coniguiente la solución general de la E.D. está dada por
		\begin{center}
			\color{red} \underline{\color{black} \(y(x) = c_1e^x \cos 2x+c_2e^x \sin 2x\)}\color{black}, \(c_1,c_2\) constantes abritrarias.
		\end{center}
	\end{alertblock}
\end{frame}

\end{document}
